\centerline{\bf \Huge Метод мат. индукции}
\centerline{\bf \Huge Знакомство}

\begin{flushright}
        \scriptsize{2...3...5...7...11...13...17...19.... Спокойствие...\\ Считай простые числа, чтобы сохранить самообладание...\\ Простые числа - это целые числа, которые делятся только на 1 и на самих себя...\\ они дают мне силу...\\Enrico Pucci. JoJo Stone ocean. Chapter 49, part 2.}
\end{flushright}

Туториал по шаганию по бесконечной лестнице:

1) \textbf{База.} Найдите начало лестницы. Первым делом нужно найти первую ступеньку и уверенно встать на неё.

2) \textbf{Переход.} После того как встали на первую ступеньку, нужно научиться перешагивать на следующую ступеньку. 

Повторяйте шаг номер 2 бесконечное количество раз. Поздравляю, вы поняли, что такое математическая индукция!

\centerline{\bf \large Для разбора}

\problem{0.0}
<<Теорема>>. При любом натуральном $n$ число $2 n+1$ четное.
<<Доказательство>>. Пусть эта <<теорема>> верна при $n=k$, то есть число $2 k+1$ четное. Докажем, что тогда число $2(k+1)+1$ также четно. Действительно, $2(k+1)+1=(2 k+1)+2$. По предположению число $2 k+1$ четно, а поэтому его сумма с четным числом 2 также четна. <<Теорема>> доказана.

\problem{0.1}
Докажите тождество: $1^2+3^2+\ldots+(2 n-1)^2=\dfrac{n\left(4 n^2-1\right)}{3}$;


\problem{0.2}
Докажите по индукции, что при любом натуральном $n$ $7^n+3 n-1$ кратно 9.

\problem{0.3}
При каких значениях $n \in \mathbb{N}$ верно неравенство $2^n>2 n+1$ ? Докажите.

\problem{0.4}
Несколько прямых делят плоскость на части. Докажите, что эти части можно раскрасить в $2$ цвета так, что граничащие части будут иметь разный цвет.

\centerline{\bf \large Формулки}

\problem{1}
Докажите тождество: $1+3+5+\ldots+(2 n-1)=n^2$ для любого натурального $n$.

\problem{2}
Докажите тождество: $1^3+2^3+\ldots+n^3=(1+2+\ldots+n)^2$ для любого натурального $n$.

\newpage

\centerline{\bf \Huge}



\problem{3}
Докажите тождество: $1 \cdot 1!+2 \cdot 2!+3 \cdot 3!+\ldots+n \cdot n!=(n+1)!-1$ для любого натурального $n$.

\centerline{\bf \large ТЧшка}

\problem{1}
Докажите по индукции, что при любом натуральном $n$ $5 \cdot 2^{3 n-2}+3^{3 n-1}$ кратно 19.

\problem{2}
Докажите, что число $11 \ldots 1$ ($3^n$ единиц) делится на $3^n$.

\problem{3}
Докажите, что существует 100---значное число, записываемое только двойками и единицами и делящееся на $2^{100}$.
%https://problems.ru/view_problem_details_new.php?id=73648

\begin{flushright}
    \vspace{-3mm}
    \textit{\scriptsize Подсказка: попробуйте заменить $100$ на $n$}
\end{flushright}

\centerline{\bf \large Неравенства}

\problem{1}
Докажите, что при любом натуральном $n$ справедливо неравенство

$$
1+\dfrac{1}{2}+\dfrac{1}{4}+\ldots+\dfrac{1}{2^n}<2
$$

\problem{2}
Докажите неравенство для натуральных $n: \dfrac{1}{\sqrt{1}}+\dfrac{1}{\sqrt{2}}+\ldots+\dfrac{1}{\sqrt{n}} \geqslant \sqrt{n}$.
%https://problems.ru/view_problem_details_new.php?id=60302

\problem{3}
Докажите неравенство для натуральных $n>1: \dfrac{1}{n+1}+\dfrac{1}{n+2}+\ldots+\dfrac{1}{2 n}>\dfrac{13}{24}$.
%https://problems.ru/view_problem_details_new.php?id=60304

\centerline{\bf \Huge}


\centerline{\bf \large Комбигеома}

\problem{1}
Докажите, что квадрат можно разрезать на $n$ (не обязательно одинаковых) квадратов для любого $n$, начиная с шести.

\problem{2}
Вершины выпуклого многоугольника раскрашены в три цвета так, что каждый цвет присутствует и никакие две соседние вершины не окрашены в один цвет. Докажите, что многоугольник можно разбить диагоналями на треугольники так, чтобы у каждого треугольника вершины были трёх разных цветов.

\problem{3}
Докажите, что для любого числа точек плоскости найдется выпуклый многоугольник с вершинами в некоторых из них, содержащий внутри себя все остальные точки.
